\newcommand{\ModuleID}{EL-II.9}
\newcommand{\ModuleTitle}{Conditions}
\newcommand{\ModuleStage}{Stage II — Core Grammar}
\newcommand{\ModuleUnit}{II.9}
\newcommand{\ModuleOutcome}{Classify and transform open, future, and contrary-to-fact conditions while preserving time, stance, and independent verb parses}
\newcommand{\ModulePrerequisites}{EL-II.8 mastery: choose conjunction-clause mood from viewpoint and relation}
\newcommand{\ModuleSessionOne}{Build the condition timeline/stance chart and parse protasis and apodosis independently, then complete Worksheet II.33 as the guided application.}
\newcommand{\ModuleSessionTwo}{Transform one lexical pair through the lesson's condition patterns, preserve reality and time in English, and complete Worksheet II.36 independently.}
\newcommand{\ModulePrintedPractice}{Worksheets II.33 (guided Variant A) and II.36 (independent Variant D), with terminal keys}
\newcommand{\ModuleExtensionPractice}{Worksheets II.34 and II.35 remain optional remediation or delayed-recall variants}
\newcommand{\ModuleNext}{EL-II.10, Relative Clauses of Characteristic}
\newcommand{\ModuleMastery}{Classify ten unseen conditions and parse every verb with no more than one error, transform one lexical pair through all five taught patterns without the table, and translate two contrary-to-fact conditions with correct English time.}
\newcommand{\ModuleDelayNote}{}
\newcommand{\ModuleReferenceFocus}{%
  \begin{itemize}
    \item protasis and apodosis as independently parsed clauses;
    \item open, future more vivid, future less vivid, present contrary-to-fact, and past contrary-to-fact patterns;
    \item \Latin{si, nisi, sin}, and conditional relatives; and
    \item English renderings that preserve possibility, expectation, or unreality and time.
  \end{itemize}}
\newcommand{\ModuleContrast}{Tense and mood must be parsed independently in both clauses before the condition is classified. Subjunctive does not by itself mean past, and indicative does not by itself promise fulfillment.}
\newcommand{\ModuleLessonSource}{\CourseRoot/shared/content/volume2/all}
\newcommand{\ModuleExerciseSource}{\CourseRoot/shared/exercises/volume2/all}
\newcommand{\ModuleWorkedBank}{II}
\newcommand{\ModuleExerciseBank}{II}
\newcommand{\ModuleWorksheetVariants}{A,D}
\newcommand{\ModuleScopeNote}{This packet teaches the five operational patterns in the course; mixed conditions require their actual clause forms and context rather than forced normalization.}
\newcommand{\ModuleEnrichment}{%
  \SubmoduleExtension
    {The conditional system}
    {A timeline/stance chart keeps chronological location separate from the speaker's presentation of likelihood or unreality.}
    {For each taught condition type, enter protasis mood/tense, apodosis mood/tense, time reference, and stance in separate columns. Add no English until the Latin columns are complete.}
    {Name the two contrary-to-fact time frames and their ordinary subjunctive tenses.}
    {Present contrary-to-fact ordinarily uses imperfect subjunctive; past contrary-to-fact uses pluperfect subjunctive. The complete chart must also preserve the indicative and future patterns taught rather than treating every condition as subjunctive.}
  \SubmoduleExtension
    {Reading the forms in context}
    {Protasis and apodosis can contain different forms, so each clause must be parsed before the pair receives a label.}
    {Take the lesson's unseen contrasts and write a full parse under every finite verb. Only then classify the condition and underline the feature that distinguishes it from its nearest neighbor.}
    {Explain why an English \English{would} cannot replace the Latin parse.}
    {The classification must follow actual tense and mood in both clauses and the time relation. English auxiliaries are translation choices and can blur distinctions; they do not establish the Latin morphology.}
  \SubmoduleExtension
    {Writing laboratory}
    {Holding vocabulary constant while changing the condition tests the grammar of stance rather than lexical recognition.}
    {Perform the lesson's five-pattern transformation. Plot each pair on the timeline/stance chart and give a close English rendering that makes the intended time and reality status audible.}
    {State which parts of the proposition remain fixed across all five versions.}
    {The lexical proposition and logical if--then relation remain fixed; verb forms, time framing, and speaker stance change according to the selected pattern. Each rendering must preserve those changes without inventing certainty.}
}
